\begin{example}[Atomic formulas]
    \label{exa:atomicFormulasHypercube}
    %The assumption of \theref{the:sufficientHLNExpressivity} is satisfied in
    Let us consider the case of atomic formulas $\formulaofat{\selindex}{\shortcatvariables}$ defined with the coordiantes at $\shortcatindicesin$ by
    \begin{align*}
        \formulaofat{\selindex}{\indexedshortcatvariables}
        = \catindexof{\selindex} \, .
    \end{align*}
    The mean polytope in this case is the $\catorder$-dimensional hypercube (see \exaref{exa:hypercubeFaces})
    \begin{align*}
        \meansetof{\atomformulaset,\ones}
        = \fullparcube
    \end{align*}
    which is called a simple polytope, since each vertex is contained in the minimal number of $\catorder$ facets.
    %Since the cube $\fullparcube$ is a face of itself, \theref{the:sufficientHLNExpressivity} implies $\hlnmeanset|_{\hlnsetof{\atomformulaset}}=\hlnmeanset$.

    Each face is characterized by the projections onto each variable, which is either $\{0\}$, $\{1\}$ or $[0,1]$.
    The projections are represented by the tuple $\hardparam$ defined in the following way:
    \begin{itemize}
        \item We define the set $\hardlegset\subset[\atomorder]$ of variables, such that the projection onto the variable is $\{0\}$ or $\{1\}$
        \item We define to each $\selindex\in\hardlegset$ an index $\headindexof{\selindex}=0$ if the projection is $\{0\}$ and $\headindexof{\selindex}=1$ if the projection is $\{1\}$.
    \end{itemize}
    Trivially, each face of the hypercube is a cube face. %$\elmeansetof{\atomformulaset}=\meansetof{\atomformulaset}$.

    % Extension: Discussion of those statistics having the hypercube as mean polytope
    More general, the mean polytope to a collection $\formulaofat{\selindex}{\shortcatvariables}$ of propositional formulas is the hypercube, if and only if for any boolean vector $\imelementwith$ we have
    \begin{align*}
        \contraction{\bigwedge_{\selindexin}\lnot^{1-\imelementat{\indexedselvariable}}\formulaofat{\selindex}{\shortcatvariables}}
        \neq 0 \, .
    \end{align*}
    This means that no conjunction of possibly negated formulas from this statistic entails nor contradicts another formula.

\end{example}